\resumo{Resumo}
O material particulado fino (PM2,5) é um poluente atmosférico crítico, associado a agravos à saúde pública, sobrecarga de sistemas de saúde e impactos ambientais. A previsão horária de curto prazo é, portanto, uma ferramenta relevante para análise ambiental, apoio a alertas e gestão da qualidade do ar. Contudo, essa tarefa é desafiadora devido à natureza não linear, ruidosa e frequentemente incompleta das séries de monitoramento, além da dependência de variáveis meteorológicas, temporais e de emissão. Este trabalho propõe um estudo experimental comparativo para previsão horária de PM2,5 com horizonte de 24 horas na estação Sapo, pertencente ao recorte CMD (Conceição do Mato Dentro, MG) do conjunto de dados. A estação foi selecionada por combinar alta cobertura de PM2,5, presença pareada de PM10 e ocorrência suficiente de eventos de maior concentração para sustentar a comparação entre modelos. O protocolo utiliza as 48 horas antecedentes como janela de entrada e inclui variáveis preditoras derivadas de PM2,5, PM10, indicadores temporais, termos de Fourier para sazonalidade e substitutos causais de PM10 disponíveis no instante da previsão. Também foi aplicado mascaramento para excluir valores imputados do alvo da função de perda e das métricas de avaliação. Foram comparados modelos LSTM, arquiteturas Seq2Seq com e sem atenção e XGBoost multissaída como modelo de referência tabular, com ajuste de hiperparâmetros por validação temporal \textit{walk-forward} e análise de robustez por múltiplas sementes. Nos resultados principais, a LSTM direta destacou-se como o modelo mais equilibrado em erro médio. O XGBoost consolidou-se como modelo de referência tabular forte, enquanto o Seq2Seq com atenção melhorou parte do comportamento em picos sem superar a LSTM direta no erro médio. Experimentos diagnósticos com variáveis auxiliares futuras elevaram o \(R^2\), evidenciando que parte da variância horária depende de informações não disponíveis causalmente no cenário operacional. Conclui-se que a modelagem deve equilibrar minimização do erro médio, preservação da amplitude dos picos e disponibilidade causal de variáveis explicativas.

\noindent \ \textbf{Palavras-chave:} PM2,5, Previsão de Séries Temporais, Aprendizado Profundo, LSTM, Seq2Seq, XGBoost, Qualidade do Ar.

\resumo{Abstract}
Fine particulate matter (PM2.5) is a critical air pollutant associated with public health burdens and environmental impacts. Short-term hourly forecasting is therefore relevant for environmental analysis, warning support, and air-quality management. However, this task is challenging because monitoring series are nonlinear, noisy, often incomplete, and influenced by meteorological, temporal, and emission-related factors. This work proposes a comparative experimental study for 24-hour-ahead hourly PM2.5 forecasting at the Sapo station, within the CMD subset (Conceição do Mato Dentro, Minas Gerais, Brazil) of the dataset. The station was selected because it combines high PM2.5 coverage, paired PM10 availability, and enough higher-concentration events to support model comparison. The protocol uses the previous 48 hours as input and includes predictors derived from PM2.5, PM10, temporal indicators, Fourier seasonality terms, and causal PM10 proxies available at prediction time. Target masking is also applied to exclude imputed target values from the loss function and evaluation metrics. The study compares LSTM models, Seq2Seq architectures with and without attention, and multi-output XGBoost as a tabular reference model, using walk-forward hyperparameter tuning and multi-seed robustness analysis. In the main results, direct LSTM stood out as the most balanced model in terms of average error. XGBoost remained a strong tabular reference model, while attention-based Seq2Seq improved part of the peak behavior without surpassing direct LSTM in average error. Diagnostic experiments with future auxiliary variables increased \(R^2\), indicating that part of the hourly variance depends on information that is not causally available in the operational setting. The results show that forecasting design must balance average-error minimization, peak-amplitude preservation, and the causal availability of explanatory variables.

\noindent \ \textbf{Keywords:} PM2.5, Time Series Forecasting, Deep Learning, LSTM, Seq2Seq, XGBoost, Air Quality.
